% =====================================================================
\chapter{The Temporal Fidelity Layer}
\label{ch:design}
% =====================================================================

This chapter answers RQ2 by design. \Cref{sec:design:goals} derives
design goals from the baseline characterisation;
\cref{sec:design:architecture} presents the architecture;
\cref{sec:design:components} details the three components; and
\cref{sec:design:integration} shows how TFL integrates into an
unmodified Renode installation.

\section{Design Goals}
\label{sec:design:goals}

\begin{description}
  \item[G1 --- Non-invasiveness.] No modification of the Renode code
        base and no recompilation: TFL must load into stock
        Renode~1.15.3 at runtime.
  \item[G2 --- Firmware transparency.] The firmware logic under test
        must not change. (The probe firmware writes to optional TFL
        registers for \emph{monitoring}, but these writes are
        no-ops on hardware-only builds and do not alter control flow.)
  \item[G3 --- Statistical realism.] TFL targets distribution-level
        fidelity: emulated timing variance should match hardware
        variance ($J \to 1$ in \cref{eq:jitterratio}), not individual
        event times.
  \item[G4 --- Configurability and reproducibility.] All temporal
        parameters (drift, jitter, delay) are set declaratively in the
        platform description; randomness draws from Renode's seeded
        PRNG so that runs remain reproducible.
  \item[G5 --- Composability.] Each deviation source is modelled by an
        independent component; any subset can be enabled for ablation
        studies.
\end{description}

\section{Architecture}
\label{sec:design:architecture}

\begin{figure}[tb]
  \centering
  \begin{tikzpicture}[
      font=\small,
      box/.style={draw, rounded corners=2pt, align=center,
                  minimum height=9mm, inner sep=4pt},
      tfl/.style={box, fill=green!12},
      std/.style={box, fill=black!6},
      arr/.style={-{Stealth[length=2.5mm]}, thick},
      lbl/.style={font=\footnotesize\itshape}
    ]
    % firmware
    \node[std, minimum width=118mm] (fw)
      {Unmodified firmware ELF (Zephyr application)};

    % system bus
    \node[std, minimum width=118mm, below=9mm of fw] (bus)
      {Renode system bus (\texttt{sysbus})};

    \draw[arr, <->] (fw) -- node[lbl, right] {load/store, IRQ} (bus);

    % standard peripherals
    \node[std, below=9mm of bus.south west, anchor=north west,
          minimum width=26mm] (cpu) {CPU\\ Cortex-M3};
    \node[std, right=4mm of cpu, minimum width=26mm] (per)
      {Stock peripherals\\ USART, GPIO, \dots};

    % TFL peripherals
    \node[tfl, right=8mm of per, minimum width=17mm] (clk)
      {\texttt{TFL\_Clock}\\ \texttt{Model}\\[1pt] {\footnotesize\srcS{3}}};
    \node[tfl, right=4mm of clk, minimum width=17mm] (mon)
      {\texttt{TFL\_Deadline}\\ \texttt{Monitor}\\[1pt] {\footnotesize\srcS{4}}};
    \node[tfl, right=4mm of mon, minimum width=17mm] (busd)
      {\texttt{TFL\_BusDelay}\\ \texttt{Model}\\[1pt] {\footnotesize\srcS{5}}};

    \draw[arr, <->] (cpu.north) -- ($(cpu.north|-bus.south)$);
    \draw[arr, <->] (per.north) -- ($(per.north|-bus.south)$);
    \draw[arr, <->] (clk.north) -- ($(clk.north|-bus.south)$);
    \draw[arr, <->] (mon.north) -- ($(mon.north|-bus.south)$);
    \draw[arr, <->] (busd.north) -- ($(busd.north|-bus.south)$);

    % TFL brace
    \node[draw, dashed, rounded corners=2pt, fit=(clk)(mon)(busd),
          inner sep=3mm, label={[lbl]below:Temporal Fidelity Layer
          (runtime-loaded C\# peripherals)}] {};

    % csv output
    \node[std, below=13mm of mon, minimum width=30mm] (csv)
      {deviation CSV\\ (per-wakeup samples)};
    \draw[arr] (mon.south) -- (csv.north);
  \end{tikzpicture}
  \caption[TFL architecture]{TFL architecture. The three TFL components
    are ordinary memory-mapped peripherals attached to the system bus by
    the platform description; Renode itself is unmodified.}
  \label{fig:tfl-arch}
\end{figure}

\Cref{fig:tfl-arch} shows the architecture. Each TFL component is a C\#
class deriving from Renode's \texttt{BasicDoubleWordPeripheral},
registered on the system bus by the \texttt{.repl} platform file. Two
integration idioms are used:

\begin{itemize}
  \item \textbf{Shadowing}: \texttt{TFL\_ClockModel} is mapped over the
        SysTick register window (\texttt{0xE000E010}--\texttt{0xE000E02F}),
        transparently perturbing every counter read the firmware
        performs. No firmware cooperation is needed.
  \item \textbf{Side-band}: \texttt{TFL\_DeadlineMonitor} and
        \texttt{TFL\_BusDelayModel} occupy otherwise-unused address
        ranges (\texttt{0x50000000}, \texttt{0x50001000}). Firmware may
        write to them for monitoring (a single \texttt{TFL\_MARK()}
        store, compiled to one \texttt{str} instruction), which is
        harmless on hardware where the region is unmapped\footnote{On
        Cortex-M3 the region is write-ignored in the probe
        configuration; guarding the macro with a build flag is equally
        possible.}.
\end{itemize}

\section{Components}
\label{sec:design:components}

\subsection{TFL\_ClockModel: Drift and Jitter Injection (\srcS{3})}
\label{sec:design:clock}

The clock model exposes a SysTick-compatible register block
(\texttt{CTRL}, \texttt{LOAD}, \texttt{CURRENT}, \texttt{CALIB}). On
every read of \texttt{CURRENT} it computes the ideal cycle count
$c_{\mathrm{true}}$ from virtual time and the configured CPU frequency
$f$, then returns a perturbed value
%
\begin{equation}
  c \;=\; c_{\mathrm{true}}
        \;+\; \underbrace{c_{\mathrm{true}}\cdot
              \frac{D}{10^6}}_{\text{drift, } D~\mathrm{ppm}}
        \;+\; \underbrace{z \cdot \sigma_J \cdot
              \frac{f}{10^9}}_{\text{jitter, } z\sim\mathcal{N}(0,1)},
  \label{eq:clockmodel}
\end{equation}
%
where $D$ (\texttt{DriftPpm}) models static oscillator offset and
$\sigma_J$ (\texttt{JitterNs}) models cycle-to-cycle noise, generated
by a Box--Muller transform from Renode's seeded PRNG (goal G4). The
evaluation configuration uses $D = 25$\,ppm and
$\sigma_J = \SI{150}{\nano\second}$, values within the tolerance band
of the STM32F103's HSE crystal circuit~\cite{stm32rm0008}.

\subsection{TFL\_DeadlineMonitor: Wakeup Deviation Tracking (\srcS{4})}
\label{sec:design:monitor}

The deadline monitor quantifies scheduler-induced deviation
\emph{inside} the emulation. Firmware writes to its \texttt{MARK}
register at every periodic wakeup; the peripheral timestamps the write
in virtual time, computes the deviation from the configured
\texttt{ExpectedPeriodNs}, updates the \texttt{LAST\_DEV\_HI/LO}
registers (readable by the firmware itself, enabling self-adaptive
tests), and streams per-sample records to a CSV file for offline
analysis. This turns the emulator into its own measurement instrument
and provides the virtual-side ground truth used in
\cref{sec:eval:tfl}.

\subsection{TFL\_BusDelayModel: Propagation Delay Injection (\srcS{5})}
\label{sec:design:bus}

The bus delay model provides a \texttt{SEND}/\texttt{RECV} register
pair backed by a bounded delay queue. A value written to \texttt{SEND}
becomes visible at \texttt{RECV} only after
%
\begin{equation}
  t_{\mathrm{deliver}} \;=\; t_{\mathrm{send}} + d + z\cdot\sigma_A,
  \qquad z\sim\mathcal{N}(0,1),
\end{equation}
%
where $d$ (\texttt{PropagationDelayNs}) is the deterministic
propagation delay and $\sigma_A$ (\texttt{ArbitrationJitterNs}) models
arbitration variance. The default $d = \SI{86806}{\nano\second}$ equals
one UART frame (10 bits) at \SI{115200}{\baud}, the wire latency that
stock Renode collapses to zero. \texttt{STATUS} exposes queue depth,
and overflow beyond \texttt{MaxQueueDepth} drops messages and counts
them in \texttt{STATS\_DROP}---reproducing the lossy behaviour of real
buffered links.

\section{Integration and Usage}
\label{sec:design:integration}

TFL is enabled entirely declaratively.
\Cref{lst:repl} shows the complete TFL platform description used in the
evaluation; \cref{lst:resc} shows the corresponding Renode script. An
ablation configuration (clock-only, deadline-only, bus-only) is a
three-line edit, satisfying goal G5.

\begin{lstlisting}[language=repl, caption={TFL platform description
  (\texttt{platforms/nucleo\_f103rb\_tfl.repl}). The base STM32F103
  description is included unchanged; TFL components are added at
  dedicated addresses.}, label={lst:repl}]
using "platforms/cpus/stm32f103.repl"

tfl_clock: Timers.TFL_ClockModel @ sysbus <0xE000E010, 0xE000E02F>
    frequency: 72000000
    DriftPpm:  25.0
    JitterNs:  150.0

tfl_monitor: Miscellaneous.TFL_DeadlineMonitor @ sysbus <0x50000000, 0x5000000F>
    ExpectedPeriodNs: 1000000
    MaxSamples: 10000

tfl_bus: Miscellaneous.TFL_BusDelayModel @ sysbus <0x50001000, 0x5000101F>
    PropagationDelayNs: 86806
    ArbitrationJitterNs: 500.0
    MaxQueueDepth: 64
\end{lstlisting}

\begin{lstlisting}[language=resc, caption={Renode execution script
  (\texttt{scripts/tfl\_full.resc}). The \texttt{i} directive compiles
  and loads the TFL C\# sources at runtime---Renode itself is never
  rebuilt.}, label={lst:resc}]
i @peripherals/TFL_ClockModel.cs
i @peripherals/TFL_DeadlineMonitor.cs
i @peripherals/TFL_BusDelayModel.cs
mach create "nucleo_f103rb_tfl_full"
machine LoadPlatformDescription @platforms/nucleo_f103rb_tfl.repl
sysbus.cpu PerformanceInMips 72
sysbus LoadELF @firmware/binaries/timing_probe_nucleo_f103rb.elf
start
\end{lstlisting}

\paragraph{Limitations by design.} TFL does not model
microarchitectural instruction timing (\srcS{1}) and does not change
how Renode recognises interrupts (\srcS{2}); the quantum remains the
lower bound for event granularity. \Cref{sec:eval:quantum} shows how
the residual \srcS{2}/\srcS{4} error scales with the quantum, and
\cref{sec:eval:discussion} discusses the consequences.
